\documentclass[12 pt, letterpaper]{hmcpset}
\usepackage{hw-preamble}

\name{Jayden Thadani}
\class{MATH 131 --- Mathematical Analysis I}
\assignment{Homework 1}
\duedate{7th September 2024}

\begin{document}

\begin{problem}[1.]
	Let $a_{0} \in \mathbb{N}$, and $a_{1}, a_{2}, \dots \in \{ 0, \dots, 9 \}$. Prove that for any  $j \in \mathbb{N}$,
	\[
		\sup \{ a_{0}.a_{1} \cdots a_{k} \mid k \ge 0 \} = \sup \{ a_{0}.a_{1} \cdots a_{k} \mid k \ge j \}
	\]
\end{problem}

\begin{solution}
	We will show that $\alpha = a_{0}.a_{1} a_{2} \dots$ is the supremum of both sets.

	First, we show that $\alpha$ is an upper bound of both sets. Any element of either set must be of the form $a_{0}.a_{1} \cdots a_{k}$ for some $k \in \mathbb{N}$. But $a_{0}.a_{1} \cdots a_{k} \le a_{0}.a_{1} \cdots a_{k} a_{k + 1}$ and thus, $a_{0}.a_{1} \cdot a_{k} \le \alpha$. Thus, $\alpha$ is an upper bound of both sets.

	Now we will show that no $\beta < \alpha$ can be an upper bound of either set. If $\beta = b_{0}.b_{1} b_{2} \cdots < \alpha$, then there must be some first $i$ such that $\beta$'s decimal expansion up to $i$ is less than that $\alpha$'s --- at $i$, we must have $b_{0}.b_{1} \cdots b_{i} < a_{0}.a_{1} \cdots a_{i}$. If $i \ge j$, we both sets would contain $a_{0}.a_{1} \cdots a_{i} > b_{0}.b_{1} \cdots b_{i}$, and thus $a_{0}.a_{1} \cdots a_{i} > \beta$. If $i < j$ then
	 \[
		a_{0}.a_{1} \cdots a_{i} \cdots a_{j} \ge a_{0}.a_{1} \cdots a_{i} > \beta.
	\]
	Thus, no $\beta < \alpha$ can be an upper bound of either set, and thus, $\alpha$ is the supremum of both sets ---
	\[
		\sup \{ a_{0}.a_{1} \cdots a_{k} \mid k \ge 0 \} = \alpha = \sup \{ a_{0}.a_{1} \cdots a_{k} \mid k \ge j \}.
	\]
\end{solution}

\pagebreak

\begin{problem}[2.]
	If $x, y$ are non-negative real numbers, prove that $x + y = y + x$ and that $xy = yx$.
\end{problem}

\begin{solution}
	Suppose $x = x_{0}.x_{1} x_{2} \cdots$ and $y = y_{0}.y_{1} y_{2} \cdots$ with $x_{0}, y_{0} \ge 0$.

	First we will show $x + y = y + x$. By definition,
	\[
		x + y = \sup \{ x_{0}.x_{1} \cdots x_{n} + y_{0}.y_{1} \cdots y_{n} \mid n \in \mathbb{N} \}.
	\]
	Notice that $x_{0}.x_{1} \cdots x_{n} + y_{0}.y_{1} \cdots y_{n} = y_{0}.y_{1} \cdot y_{n} + x_{0}.x_{1} \cdots x_{n}$ because addition of rationals is commutative. Thus,
	 \[
		\begin{split}
			x + y & = \sup \{ x_{0}.x_{1} \cdots x_{n} + y_{0}.y_{1} \cdots y_{n} \mid n \in \mathbb{N} \} \\
			      & = \sup \{ y_{0}.y_{1} + \cdots + y_{n} + x_{0}.x_{1} + \cdots + x_{n} \mid n \in \mathbb{N} \} \\
			      & = y + x
		\end{split}
	\]
	as desired. Note that the equality of the supremums is because the sets are equal.

	Similarly, we can show that $xy = yx$. By definition,
	\[
		xy = \sup \{ (x_{0}.x_{1} \cdots x_{n})(y_{0}.y_{1} \cdots y_{n}) \mid n \in \mathbb{N} \}.
	\]
	Also $(x_{0}.x_{1} \cdots x_{n}) (y_{0}.y_{1} \cdots y_{n}) = (y_{0}.y_{1} \cdot y_{n}) (x_{0}.x_{1} \cdots x_{n})$ since multiplication is commutative in $\mathbb{Q}$. Thus,
	\[
		\begin{split}
			xy & = \sup \{ (x_{0}.x_{1} \cdots x_{n})(y_{0}.y_{1} \cdots y_{n}) \mid n \in \mathbb{N} \} \\
			      & = \sup \{ (y_{0}.y_{1} \cdot y_{n}) (x_{0}.x_{1} \cdots x_{n}) \mid n \in \mathbb{N} \} \\
			      & = yx
		\end{split}
	\]
	as desired. Again, the equality of the supremums is because the sets are equal.
\end{solution}

\pagebreak

\begin{problem}[3.]
	If $x$, $y$, and $z$ are non-negative real numbers with $x < y$, prove that $x + z < y + z$.
\end{problem}

\begin{solution}
	Suppose $x = x_{0}.x_{1} x_{2} \cdots$, $y = y_{0}.y_{1} y_{2} \cdots$, and  $z = z_{0}.z_{1} z_{2} \cdots$ where $x_{0}, y_{0}, z_{0} \ge 0$.

	If $x < y$, we know that $x_{0}.x_{1} \cdots x_{n} < y_{0}.y_{1} \cdots y_{n}$ for each $n \in \mathbb{N}$. Since the rationals are an ordered field, $x_{0}.x_{1} \cdots x_{n} + z_{0}.z_{1} \cdots z_{n} < y_{0}.y_{1} \cdots y_{n} + z_{0}.z_{1} \cdots z_{n}$ for each $n \in \mathbb{N}$. We also know that there must be some first $i$ such that $x_{0}.x_{1} \cdots x_{i} < y_{0}.y_{1} \cdots y_{i}$. In particular, this means that $y_{0}.y_{1} \cdots y_{i} - x_{0}.x_{1} \cdots x_{i} \ge 1/10^{i}$, and also $y_{0}.y_{1} \cdots y_{i} - x_{0}.x_{1} \cdots x_{n} \ge 1/10^{i}$ for any $n \in \mathbb{N}$. Finally, notice that $z_{0}.z_{1} \cdots z_{n} - z_{0}.z_{1} \cdots z_{i} < 1/10^{i + 1}$.

	Thus,
	\[
		(y_{0}.y_{1} \cdots y_{i} + z_{0}.z_{1} \cdots z_{i}) - (x_{0}.x_{1} \cdots x_{n} + z_{0}.z_{1} \cdots z_{n}) > \frac{1}{10^{i}} - \frac{1}{10^{i + 1}}.
	\]
	for any $n \in \mathbb{N}$. Notice that by rearranging we get
	\[
		y_{0}.y_{1} \cdots y_{i} + z_{0}.z_{1} \cdots z_{i} - (1/10^{i} - 1/10^{i + 1}) > x_{0}.x_{1} \cdots x_{n} + z_{0}.z_{1} \cdots z_{n}.
	\]
	That is, $y_{0}.y_{1} \cdots y_{i} + z_{0}.z_{1} \cdots z_{i} - (1/10^{i} - 1/10^{i + 1})$ is an upper bound of $\{ x_{0}.x_{1} \cdots x_{n} + z_{0}.z_{1} \cdots z_{n} \mid n \in \mathbb{N} \}$. Recall that $x + z = \sup \{ x_{0}.x_{1} \cdots x_{n} + z_{0}.z_{1} \cdots z_{n} \mid n \in \mathbb{N} \}$.

	The chain of inequality follows
	\[
		\begin{split}
			y + z & \ge y_{0}.y_{1} \cdots y_{i} + z_{0}.z_{1} \cdots z_{i} \\
			      & > y_{0}.y_{1} \cdots y_{i} + z_{0}.z_{1} \cdots z_{i} - (1/10^{i} - 1/10^{i + 1}) \\
			      & \ge x + z.
		\end{split}
	\]
	That is, $y + z > x + z$, as desired.
\end{solution}

\pagebreak

\begin{problem}[4.]
	Let $\varepsilon = a_{0}.a_{1} \cdots$ be a positive real number. Prove that there exists $k > 1$ such that $10^{k} \varepsilon > 1$.
\end{problem}

\begin{solution}
	Since $\varepsilon$ is positive, we know that there must be some first $k$ such that $a_{k - 1} > 0$ (else, $\varepsilon \le 0$). Then consider $10^{k} \varepsilon$. By our definition of multiplication we have
	\[
		10^{k} \varepsilon = \sup \{ 10^{k} (a_{0}.a_{1} \cdots a_{n}) \mid n \in \mathbb{Z}_{\ge 0} \}.
	\]
	The supremum cannot be less than any element of the set. However, we know that the set contains
	\[
		\begin{split}
			10^{k} (a_{0}.a_{1} \cdots a_{k}) & = \sum_{n = 0}^{k} 10^{k - n} a_{n} \\
							  & \ge 10 a_{k - 1} \\
							  & \ge 10 \\
							  & > 1.
		\end{split}
	\]
	Thus,
	\[
		\begin{split}
			10^{k} \varepsilon & = \sup \{ 10^{k}(a_{0}.a_{1} \cdots a_{n}) \mid n \in \mathbb{Z}_{\ge 0} \} \\
					   & \ge 10^{k} (a_{0}.a_{1} \cdots a_{k}) \\
					   & > 1.
		\end{split}
	\]
	as desired.
\end{solution}

\pagebreak

\begin{problem}[5.]
	Let $x = a_{0}.a_{1} \cdots$ and $y = b_{0}.b_{1} \cdots$ be positive real numbers. assuming that $x < y$ prove that there exists a decimal number $w$ such that $x < w < y$.
\end{problem}

\begin{solution}
	Consider $\varepsilon = y - x$. We know, by the previous exercise, that there exists some $k$ such that $10^{k} \varepsilon > 1$, and thus, $1/10^{k} < \varepsilon$. Note that $\alpha = 0.0 \cdots 0 a_{k + 1} a_{k + 2} \cdots$, the trailing decimals after $a_{k}$ give $\alpha < 1/10^{k}$. Thus, $0 < 1/10^{k} - \alpha < \varepsilon$. This gives us $x < x + 1/10^{k} - \alpha < y$.

	Now we only need to show that $x - \alpha$ is the decimal $a_{0}.a_{1} \cdots a_{k}$. If we consider. Recall that
	\[
		x - \alpha = \sup \{ a_{0}.a_{1} \cdots a_{n} - 0.0 \cdots 0 a_{k + 1} \cdots a_{n} \mid n \in \mathbb{Z}_{\ge 0} \}.
	\]
	If $n \le k$, we have the corresponding element $a_{0}.a_{1} \cdots a_{n} - 0 = a_{0}.a_{1} \cdots a_{n} \le a_{0}.a_{1} \cdots a_{k}$ in the set. If $n > k$, we have $a_{0}.a_{1} \cdots a_{k} \cdots a_{n} - 0.0 \cdots 0 a_{k + 1} \cdots a_{n} = a_{0}.a_{1} \cdots a_{k}$ in the set. Thus, $a_{0}.a_{1} \cdots a_{k}$ is an upper bound on the set, and is contained in the set. Thus, it must be the supremum. 

	Thus is, $x -  \alpha$ is a decimal, and so is $w = x - \alpha + 1/10^{k}$ which we've already shown satisfies $x < w < y$.
\end{solution}

\pagebreak

\begin{problem}[6.]
	Prove that if $p$ and $q$ are decimal numbers then there exists an irrational number $z$ such that $p < z < q$.
\end{problem}

\begin{solution}
	By the previous exercise, we know that there exists a (rational) decimal number $w$ with
	\[
		\frac{p}{\sqrt{2}} < w < \frac{q}{\sqrt{2}}.
	\]
	Thus,
	\[
		p < w  \sqrt{2} < q.
	\]
	Note that $w \sqrt{2}$ cannot be rational --- otherwise $(1 / w) w \sqrt{2} = \sqrt{2}$ would be rational (since the rationals have inverses and are closed under multiplication). Thus, $z = w \sqrt{2}$ is the desired irrational between $p$ and $q$.
\end{solution}

\end{document}
